\Needspace{20\baselineskip}
\section{The full class machine and concrete forward simulation}
\label{sec:uam-classes}
The determinism required by the abstract counting-simulation theorem is
literal equality of successors. Named hypernets instead have unique successors
up to recursive focus-preserving isomorphism. This section constructs the full
UAM on those isomorphism classes, proves its literal determinism, constructs
its concrete specimen interface, and obtains the concrete forward counting
simulation and quantitative refinement of remote rooted specimens.

The behaviour is required to respect recursive endpoint isomorphism when
using its named-state determinism hypothesis. A closure construction below
shows that this invariance concerns representative names: closing compute
rules under renaming leaves the entire class transition relation unchanged.
No deterministic transition on raw carrier names is assumed.

Safety and robustness here are evaluated directly on the full class machine,
with all the concrete clauses retained. We do not infer class safety merely
from the previously defined named-state safety: projection of finite runs
alone would not justify preservation of stuckness or refinement. General
initial observation lifting and the converse integration remain separate.

\Needspace{14\baselineskip}
\subsection{Comparing steps from isomorphic sources}
\begin{definition}[Renaming a presentation's source]
A substitution presentation can replace its source by any recursively
isomorphic state while leaving its base graph, selected variable, and target
unchanged. Compose the presentation's source isomorphism with the given one.
\end{definition}
\begin{theorem}[Substitution uniqueness across source classes]
Two substitution steps whose sources are recursively isomorphic have
recursively isomorphic targets.
\end{theorem}
\begin{proof}
Extract a presentation of each substitution. Rename the first presentation's
source to the second source, then apply the already proved uniqueness theorem
for presentations of a common source. This comparison does not assert a
new literal substitution transition at the renamed source; the presentation
is sufficient to identify the two actual reducts.
\end{proof}

\lean{Graph/UAMClassDeterminism}{10}{12}{renameSource}

\lean{Graph/UAMClassDeterminism}{17}{23}{substitution_iso_deterministic}

\begin{lemma}[Substitution and compute remain exclusive under renaming]
A substitution source cannot be recursively isomorphic to a compute source.
\end{lemma}
\begin{proof}
The substitution exposes a variable with singleton source at the focus target.
The source isomorphism transports its primitive variable label, ordered source
list, and focus target to the other graph. Computation exposes an active
operation with that same singleton source. Source uniqueness identifies the
two edges, contradicting their different primitive labels.
\end{proof}

\lean{Graph/UAMClassDeterminism}{26}{38}{substitution_compute_iso_exclusive}

\Needspace{10\baselineskip}
\termdef{behaviour-iso-invariance}
\begin{definition}[Behaviour invariance under names]
An extrinsic behaviour respects isomorphism if independently renaming either
endpoint recursively preserves every behaviour transition. Inverse
isomorphisms also give reflection. This condition is distinct from determinism,
refocusing, and locality of a compute rule.
\end{definition}

\lean{Graph/UAMClassDeterminism}{44}{46}{RespectsIso}

\begin{theorem}[Full-UAM uniqueness across isomorphic sources]
If the behaviour respects isomorphism and is deterministic modulo recursive
isomorphism at each common named source, full-UAM steps from recursively
isomorphic sources have recursively isomorphic targets.
\end{theorem}
\begin{proof}
Inspect the two transition families. For two search steps, transport the first
source recursively and use search determinism. For two substitutions, use
uniqueness across source classes. For two extrinsic steps, transport the first
source while fixing its target and apply behaviour determinism at the now
common source. Search versus rewrite is impossible because source modes are
preserved by isomorphism. Substitution versus compute is excluded by the
preceding incidence lemma, in either direction. These cases exhaust the UAM.
\end{proof}

\lean{Graph/UAMClassDeterminism}{53}{73}{iso_deterministic}

\Needspace{14\baselineskip}
\subsection{The full UAM on recursive state classes}
\Needspace{10\baselineskip}
\termdef{full-class-step}
\begin{definition}[Full class transition]
A transition from class $s$ to class $t$ exists when there are named states
$a,b$ representing those respective classes and an actual full-UAM step
$a\to b$. No global representative is selected. The original named transition
projects by taking its own endpoints as witnesses.
\end{definition}

\lean{Graph/UAMClasses}{11}{12}{Step}

\lean{Graph/UAMClasses}{15}{16}{project}

\begin{theorem}[Literal class determinism]
Under the preceding behaviour hypotheses, two full transitions from the same
class have equal target classes.
\end{theorem}
\begin{proof}
Extract the two actual source and target pairs. Equality of their source
classes yields a recursive focus-preserving source isomorphism by quotient
exactness. The full-UAM comparison theorem supplies a target isomorphism.
Quotient soundness identifies the two target classes, and their original
endpoint equalities finish the proof.
\end{proof}
\begin{lemma}[Finality and the class machine]
A final class has no full transition. The full class machine uses this
transition relation and the existing class-invariant boundary finality.
\end{lemma}
\begin{proof}
For a purported successor, extract its representative step. Transfer finality
to that source representative and apply the concrete full-UAM final-normality
theorem. This proves the machine's final-normality field.
\end{proof}

\lean{Graph/UAMClasses}{19}{24}{deterministic}

\lean{Graph/UAMClasses}{27}{30}{final_normal}

\lean{Graph/UAMClasses}{33}{36}{machine}

\begin{lemma}[Rootedness and exact run projection]
Refocusing of the extrinsic behaviour makes every full class step preserve
rootedness. Every named finite run projects to a class run of exactly the same
length. Every class run preserves rootedness.
\end{lemma}
\begin{proof}
For a single class step, use its representative equalities to transfer the
rooted premise, apply concrete full-UAM refocusing, and transfer the result
back. For run projection, induct on the named run, mapping each transition and
retaining its successor count. For the class-run invariant, induct on the run
and apply the single-step rootedness theorem. The zero-step cases are identities.
\end{proof}

\lean{Graph/UAMClasses}{39}{45}{rooted_step}

\lean{Graph/UAMClasses}{48}{52}{project_steps}

\lean{Graph/UAMClasses}{73}{78}{rooted_steps}

\begin{lemma}[Recursive search inclusion and fixed-source reflection]
Recursive search gives a full class transition. Conversely, a class transition
from the class of a specified non-rewrite state admits a recursively renamed
search step from that very state to a representative of the target class.
\end{lemma}
\begin{proof}
For inclusion, factor recursive search into literal search and a recursive
target isomorphism. Project the literal transition and identify its target
class with the required one. For reflection, extract the class transition's
actual representatives and obtain a source isomorphism by quotient exactness.
A rewrite representative would have rewrite mode and hence contradict the
specified source mode. The step is therefore search. Transport its source
along the inverse isomorphism, retaining its actual target and target-class
equality. This reflection theorem concerns non-rewrite sources; it does not
claim that arbitrary rewrite runs lift to every raw representative.
\end{proof}

\lean{Graph/UAMClasses}{55}{59}{search}

\lean{Graph/UAMClasses}{63}{70}{search_of_nonrewrite}

\Needspace{14\baselineskip}
\subsection{The class-valued specimen interface}
\Needspace{10\baselineskip}
\termdef{class-specimen-system}
\begin{definition}[Class-valued specimen system]
Retain the actual allowed recursive specimens, their mode, and their entering
predicate, but take the classes of their filled states as endpoints. Use the
existing rooted predicate on classes. Supply rooted-step preservation from
the class refocusing theorem and final agreement from the concrete boundary
law.
\end{definition}
\begin{proof}[Construction of context traversal]
For a settled specimen, either mode alternative excludes rewrite. Reflect an
available class transition from either endpoint to recursive search at its
specified actual representative. The concrete successor-specimen theorem
constructs a remote allowed specimen with synchronized search transitions.
Include these in the full class relation. This supplies the final interface
field for arbitrary source representatives, without postulating a class
locality law.
\end{proof}

\lean{Graph/ClassSpecimenSystem}{13}{34}{system}

\begin{lemma}[Concrete class lifting]
Class lifting is exactly existence of a remote concrete specimen with rooted
endpoints and the specified endpoint classes. Named lifting projects into
class lifting. The previously constructed initial relocation belongs to
class lifting.
\end{lemma}
\begin{proof}
The first statement is definitional. For projection, unpack a named lifting
witness, keep its specimen and predicates, and apply the class constructor
to its endpoint equalities. For initial relocation, both endpoints are initial
by the constructed initiality theorem, hence rooted by zero-step runs.
Its search mode makes it remote, and its endpoint-class equalities are reflexive.
This last theorem starts with an existing specimen and is not yet the general
observation-lifting theorem for arbitrary context and fragment observations.
\end{proof}

\lean{Graph/ClassSpecimenSystem}{38}{42}{lifting_iff}

\lean{Graph/ClassSpecimenSystem}{45}{50}{lifting_project}

\lean{Graph/ClassSpecimenSystem}{53}{59}{initial_lifting}

\Needspace{14\baselineskip}
\subsection{Class safety, robustness, and forward sufficiency}
\begin{definition}[Mode of a state class]
The focus mode of a class is the focus mode of any representative. Recursive
focus-preserving isomorphisms prove the required equality of modes, so quotient
lifting is well-defined.
\end{definition}

\lean{Graph/ClassSafety}{8}{9}{mode}

\Needspace{10\baselineskip}
\termdef{class-quasi-specimen}
\begin{definition}[Class quasi-specimen]
A quasi-specimen between two classes has an unconditional concrete specimen
witness whose mode agrees with both endpoint modes. Conditional on rootedness
of both classes, its filled endpoints are rooted and the two specified
comparison relations hold. The non-rewrite variant additionally excludes
rewrite mode at the first endpoint.
\end{definition}
This retains the full witness and mode clauses of the earlier concrete
predicate, with comparisons evaluated on the class machine.

\lean{Graph/ClassSafety}{18}{24}{Quasi}

\lean{Graph/ClassSafety}{27}{28}{NonRewriteQuasi}

\Needspace{10\baselineskip}
\termdef{class-input-safety}
\begin{definition}[Input safety on the class machine]
Every entering search specimen satisfies one of the same three alternatives:
both endpoint classes run to stuck states; or the two runs reach a settled
specimen with lengths $1+k,l$ satisfying the cost relation; or lengths
$1+k,1+l$ satisfy the cost relation and reach a non-rewrite quasi-specimen
up to equivalence. No source-rootedness restriction is imposed.
\end{definition}

\lean{Graph/ClassSafety}{31}{42}{InputSafe}

\Needspace{10\baselineskip}
\termdef{class-robustness}
\begin{definition}[Robustness on the class machine]
For every non-entering rewrite specimen and every successor class of its left
endpoint, either a source side is unrooted, the successor runs to a stuck
class, or the two sides have appropriately cost-related runs to a non-rewrite
quasi-specimen up to the specified refinements. Every successor class is
quantified, independently of the names of the transition witnesses.
\end{definition}
Since the source mode is rewrite and mode is invariant under isomorphism,
search cannot witness such a successor. Thus the full-step quantifier here
ranges precisely over the applicable rewrite transitions.

\lean{Graph/ClassSafety}{46}{53}{Robust}

\begin{lemma}[The three class bridges]
Class quasi-specimens yield the abstract rooted quasi predicate. Class input
safety and robustness yield the corresponding laws of the constructed
class-valued specimen system.
\end{lemma}
\begin{proof}
For quasi-specimens, receive the two rootedness hypotheses, use the witness's
conditional comparisons, and transport the endpoint's non-rewrite mode along
its equality with the witness mode. For input safety, retain the stuck and
direct alternatives unchanged; in the third alternative retain costs and runs
and apply the quasi bridge, keeping equivalence on both sides. For robustness,
the given rootedness contradicts exactly the unrooted disjunct. The stuck-tail
case is unchanged. In the matching case retain both runs and the cost comparison
and use the quasi bridge with the given refinement relations.
\end{proof}

\lean{Graph/ClassSafety}{59}{63}{quasi_bridge}

\lean{Graph/ClassSafety}{66}{73}{inputSafe_bridge}

\lean{Graph/ClassSafety}{76}{83}{robust_bridge}

\begin{theorem}[Concrete forward counting simulation]
Assume behaviour invariance, behaviour determinism modulo isomorphism,
behaviour refocusing, allowed-context invariance, output closure, the reasonable
cost-order hypotheses, and the class input-safety and robustness predicates.
Then the concrete class lifting is a counting simulation for the full class UAM.
\end{theorem}
\begin{proof}
Instantiate the abstract forward counting-simulation theorem with the constructed
class specimen system. Class determinism supplies its literal determinism
premise; the two concrete bridges supply safety and robustness. All locality
fields have already been proved from the concrete hypernet semantics.
\end{proof}
\begin{corollary}[Quantitative refinement of concrete specimens]
Under the same hypotheses, the left endpoint class of every remote jointly
rooted concrete specimen refines its right endpoint class with the given
cost order.
\end{corollary}
\begin{proof}
The specimen, remoteness, rootedness, and reflexive class equalities form a
lifting witness. Apply counting-simulation soundness using literal class
determinism and the reasonable cost orders. The remaining contextual step is
to show that arbitrary allowed initial observations give such witnesses.
\end{proof}

\lean{Graph/ClassSafety}{86}{94}{countingSimulation}

\lean{Graph/ClassSafety}{99}{106}{sufficiency}

\Needspace{14\baselineskip}
\subsection{Behaviour closure and non-vacuous regressions}
\begin{lemma}[Compute shape respects recursive renaming]
Independent recursive renamings of the two endpoints preserve compute shape.
\end{lemma}
\begin{proof}
Transport both focus modes. Carry the active-operation edge through the source
edge bijection; primitive-label equivalence preserves its operation symbol and
hence activity. Mapping its singleton source list gives the renamed focus
target, because the isomorphism preserves that target.
\end{proof}

\lean{Graph/BehaviourIsoClosure}{9}{15}{State.Compute.rename}

\Needspace{10\baselineskip}
\termdef{behaviour-iso-closure}
\begin{definition}[Behaviour closure under names]
The isomorphism closure of a behaviour contains a transition when an original
behaviour transition has recursively isomorphic endpoints. The preceding
compute-shape lemma verifies that this is again an admissible behaviour.
\end{definition}
\begin{lemma}[Inclusion and invariance]
The closure contains each original transition and respects independent
recursive endpoint renaming.
\end{lemma}
\begin{proof}
Use identity isomorphisms for inclusion. For invariance, compose each original
witness isomorphism with the corresponding new endpoint isomorphism.
\end{proof}

\lean{Graph/BehaviourIsoClosure}{20}{25}{isoClosure}

\lean{Graph/BehaviourIsoClosure}{28}{29}{into_isoClosure}

\lean{Graph/BehaviourIsoClosure}{32}{34}{isoClosure_invariant}

\begin{theorem}[Behaviour closure preserves the full class relation]
Replacing the behaviour by its isomorphism closure leaves every full-UAM
class transition unchanged.
\end{theorem}
\begin{proof}
Extract a representative step. Intrinsic steps are unchanged. For an extrinsic
closure step, take its original compute witnesses and compose the endpoint-class
equalities with quotient soundness for their isomorphisms. This gives an original
class step. Conversely, retain the representatives and include any original
extrinsic step in the closure. Thus closing under names adds no new class
computation. The theorem does not assert that an arbitrary behaviour becomes
deterministic after closure.
\end{proof}

\lean{Graph/BehaviourIsoClosure}{37}{48}{isoClosure_class_step}

The substitution regression projects its actual two-step terminating execution
to an exactly two-step class execution. The compute regression closes the earlier
nonempty singleton behaviour under recursive renaming. Every result is isomorphic
to its one fixed initial result, so two results are isomorphic to each other;
initiality also makes each result rooted. The general theorem therefore proves
literal determinism of this nonempty full class machine. Projection supplies
its genuine extrinsic transition, and the closure theorem checks equality with
the original singleton's class relation.

\lean{Graph/UAMClassExamples}{10}{15}{two_step_execution}

\lean{Graph/UAMClassExamples}{18}{18}{behaviour}

\lean{Graph/UAMClassExamples}{21}{23}{behaviour_deterministic}

\lean{Graph/UAMClassExamples}{26}{28}{behaviour_refocusing}

\lean{Graph/UAMClassExamples}{31}{33}{deterministic}

\lean{Graph/UAMClassExamples}{36}{38}{compute_step}

\lean{Graph/UAMClassExamples}{41}{43}{same_class_steps}

Finally the nested-hole example instantiates the class-valued specimen system
with empty behaviour and the previously proved output-closed template. Its
constructed initial relocation is a concrete class-lifting witness. This checks
that the interface is exercised on a real recursive specimen, not only an empty
specimen type.

\lean{Graph/UAMClassExamples}{46}{48}{specimenSystem}

\lean{Graph/UAMClassExamples}{50}{55}{initial_lifting}

\paragraph{Status.}
The determinism mismatch, full class machine, class locality, and forward
concrete counting simulation are now discharged. General observation lifting
and the converse safety/robustness integration remain, followed by the complete
contextual theorem. The named-state and class safety predicates are kept
explicitly distinct; no unproved equivalence between them is used. The planning
estimate is now about 93\% of mathematical effort, with a plausible range of
80--96\%. Appendix~\ref{app:remaining-work} retains only unfinished obligations.
